\section{From the algebraic certificate to a geometric flux}
\label{flux:section}
\label{sec:flux}

We work on the regular part of the spherical component $\Omega\subset S^5$
constructed in Section~\ref{sec:reductions}, and put
\[
  \Omega_+=\{x\in\Omega:w(x)>0\},\qquad a=\frac{29}{50}.
\]
All assertions up to the side boundary use derivatives from the positive
side.  In the suspended case the two tips are excluded throughout
Section~\ref{sec:flux}.  We apply the universal inequalities of
Proposition~\ref{cert:universal}, whose proof follows in
Section~\ref{sec:certificate-proof}.  For its coefficient polynomials
$f,U,K,M,P$, with their stated slot orders, set
\begin{equation}\label{flux:CG}
 C_{lijk}=K_{lijk}-K_{iljk},\qquad
 G_{lijk}=\delta_{li}U_{jk}+C_{lijk}.
\end{equation}
Thus $C_{iljk}=-C_{lijk}$.  No other slot symmetry of $U$ or $K$ will be
needed.  The orthogonal equivariance asserted in
Proposition~\ref{cert:universal} ensures that all the fields defined below
are independent of the orthonormal frame.

\subsection{Normalized differentiation and the fourth derivatives}

On $\Omega_+$ define
\begin{equation}\label{flux:normalized-data}
 A=\frac{T}{w},\qquad X=\frac{B}{w},\qquad
 q_l=X_{ljk}A_{jk}.
\end{equation}
The identities of Proposition~\ref{geo:interior-identities} show that
$A\in\STF^2(\R^5)$, $|A|=1$, and $X\in\STF^3(\R^5)$.
Moreover, $g^2+|p|^2\leq1$, so these are admissible data for the
interior certificate.  The homogenizing variable in that certificate is
set to $z=1$ in every geometric application.

\begin{lemma}\label{flux:normalized-differentiation}
At the data \eqref{flux:normalized-data}, the derivation $D_l$ of
\eqref{cert:D}, evaluated at $z=1$, equals covariant differentiation on
every coefficient polynomial in $A,p,g,\delta$.  In particular,
\begin{equation}\label{flux:first-derivatives}
 \nabla_lw=wq_l,\qquad
 \nabla_lA_{ij}=X_{lij}-q_lA_{ij},\qquad
 \nabla_lp_i=wA_{li}-g\delta_{li},\qquad
 \nabla_lg=p_l.
\end{equation}
Furthermore,
\begin{equation}\label{flux:fourth-elimination}
 \frac1wG_{lijk}\nabla_lB_{ijk}
 =5U_{jk}A_{jk}
   +C_{lijk}(\delta_{lj}A_{ik}+\delta_{lk}A_{ij}).
\end{equation}
\end{lemma}

\begin{proof}
Differentiating $w^2=T_{jk}T_{jk}$ gives
$\nabla_lw=T_{jk}B_{ljk}/w=wq_l$.  The quotient rule then gives the
formula for $\nabla A$, and $T=\nabla^2g+g\Id$ gives the formula for
$\nabla p$.  In a covariantly constant orthonormal frame at the point
under consideration these are exactly the generator rules for $D_l$.
The product rule and compatibility of the connection with contractions
prove the assertion for every coefficient polynomial.  The rules also
preserve the normalization constraints:
\[
 D_l(\tr A)=0,\qquad D_l(|A|^2)=2q_l-2q_l=0
\]
on the admissible data.

Here $D_l$ is applied only to coefficient polynomials in
$A,p,g,\delta$.  Differentiation of the factors $w$ and $X$, which are
not coefficient variables, is performed separately.  In particular, no
rule for $D_lX$ or $D_lq$ is being assumed.

For the $U$ part of \eqref{flux:fourth-elimination},
\[
 \frac1w\delta_{li}U_{jk}\nabla_lB_{ijk}
 =\frac1wU_{jk}\Delta T_{jk}=5U_{jk}A_{jk}.
\]
For the skew part, the curvature identity in
Proposition~\ref{geo:interior-identities} yields
\begin{align*}
 \frac1w C_{lijk}\nabla_lB_{ijk}
 &=\frac1{2w}C_{lijk}
       (\nabla_lB_{ijk}-\nabla_iB_{ljk})\\
 &=\frac12 C_{lijk}
   (\delta_{lj}A_{ik}+\delta_{lk}A_{ij}
       -\delta_{ij}A_{lk}-\delta_{ik}A_{lj})\\
 &=C_{lijk}(\delta_{lj}A_{ik}+\delta_{lk}A_{ij}).
\end{align*}
The last equality follows by interchanging the summed indices $l,i$
in the last two terms and using $C_{iljk}=-C_{lijk}$.  Equivalently,
$C_{lijk}\nabla_lB_{ijk}$ equals
$K_{lijk}(\nabla_lB_{ijk}-\nabla_iB_{ljk})$; thus the factor $1/2$
belongs to the formula with $C$, not to the formula with $K$.

This calculation eliminates every undetermined fourth derivative of
$g$.  Indeed, if two possible tensors $\nabla_lB_{ijk}$ have the same
Laplacian contraction and the same skew difference in $l,i$, their
difference has zero contraction with $\delta_{li}U_{jk}$ and with
$C_{lijk}$.  No constraint on the remaining fourth derivative entries is
required.
\end{proof}

\subsection{The square root and the interior comparison}

Proposition~\ref{cert:universal} gives
\begin{equation}\label{flux:f-positive}
 f(A,p,g)\geq\epsilon_*
   =\frac{438907}{12000000}>0
\end{equation}
on the entire normalized coefficient domain.  The verification in
Section~\ref{cert:interior-remainder} obtains this bound using the free
algebraic data
$X=w=z=0$.  Since $f$ is independent of $w$, this permitted algebraic
substitution does not impose $w=0$ on the geometric data: the fields
below are defined on $\Omega_+$, where $w>0$.

Define on $\Omega_+$
\begin{equation}\label{flux:fields}
 h=w^a\sqrt f,\qquad
 Z_l=w^{2a}\bigl(G_{lijk}X_{ijk}+wM_l+P_l\bigr).
\end{equation}
These fields are smooth on $\Omega_+$ and extend smoothly to each
regular side boundary point, since $w>0$ there by
Proposition~\ref{geo:boundary-jets}.

\begin{proposition}\label{flux:interior-comparison}
The scalar $h$ and vector field $Z$ satisfy
\begin{equation}\label{flux:interior-inequality}
 \diver Z\geq 4h^2+|\nabla h|^2+\epsilon_*w^{2a}
 \qquad\text{on }\Omega_+.
\end{equation}
\end{proposition}

\begin{proof}
First,
\[
 \nabla_lX_{ijk}=\frac1w\nabla_lB_{ijk}-q_lX_{ijk}.
\]
Applying the product rule to each of the three summands of $Z$, and
then using \eqref{flux:fourth-elimination}, gives the exact identity
\begin{align}
 w^{-2a}\diver Z
 ={}&5U_{jk}A_{jk}
       +C_{lijk}(\delta_{lj}A_{ik}+\delta_{lk}A_{ij})\notag\\
 &+\bigl[D_lG_{lijk}+(2a-1)q_lG_{lijk}\bigr]X_{ijk}\notag\\
 &+w\bigl[D_lM_l+(2a+1)q_lM_l\bigr]
       +D_lP_l+2aq_lP_l.\label{flux:divergence-identity}
\end{align}
Every $D_l$ in Equation~\eqref{flux:divergence-identity} is evaluated at $z=1$.
The three coefficients of $q$ come respectively from differentiating
$w^{2a}X$, $w^{2a+1}$, and $w^{2a}$; in the first expression the quotient
$X=B/w$ contributes $-q_lX$.  Thus their values are exactly
\[
 2a-1=\frac4{25},\qquad 2a+1=\frac{54}{25},\qquad
 2a=\frac{29}{25}.
\]

Put
\[
 \mathcal F_l=D_lf+2aq_lf,\qquad
 \ell_* =\frac{\mathcal F}{2f}.
\]
Differentiating $h$ and completing the square give, respectively,
\begin{equation}\label{flux:square-root-derivative}
 \nabla_lh=\frac{w^a}{2\sqrt f}\mathcal F_l,
 \qquad
 |\nabla h|^2=\frac{w^{2a}}{4f}|\mathcal F|^2,
\end{equation}
and
\[
 f|\ell|^2-\mathcal F\cdot\ell
 =f|\ell-\ell_*|^2-\frac{|\mathcal F|^2}{4f}.
\]
In particular, substituting \eqref{flux:divergence-identity} into the
interior polynomial $\mathcal I$ of \eqref{cert:interior} gives the
identity, for every $\ell$,
\begin{equation}\label{flux:square-completion}
 \mathcal I\big|_{z=1}
 =w^{-2a}\bigl(\diver Z-4h^2-|\nabla h|^2\bigr)
      +f|\ell-\ell_*|^2.
\end{equation}
The division by $f$ is justified uniformly by
\eqref{flux:f-positive}.  Evaluating \eqref{cert:interior} at
$\ell=\ell_*$ therefore gives the slightly stronger estimate
\[
 \diver Z-4h^2-|\nabla h|^2
 \geq\epsilon_*w^{2a}(1+|\ell_*|^2),
\]
which implies \eqref{flux:interior-inequality}.
\end{proof}

\subsection{Boundary reflection and the signed flux}

At a regular side boundary point, Proposition~\ref{geo:boundary-jets} gives
\[
 g=0,\quad p=N,\quad |p|=1,\quad Ap=\lambda p,
 \qquad \lambda=-\frac H w<0.
\]
As in the boundary certificate, define
\begin{equation}\label{flux:boundary-model}
 S_0=\lambda A-A^2,\qquad
 V_{ijk}=p_i(S_0)_{jk}+p_j(S_0)_{ki}+p_k(S_0)_{ij}
          +p_ip_jp_k.
\end{equation}

\begin{lemma}\label{flux:boundary-parity}
At such a point,
\begin{equation}\label{flux:boundary-contraction}
 p_lG_{lijk}\frac{B_{ijk}}{w^2}
 =p_lG_{lijk}V_{ijk}
 =U_{jk}(S_0+p\otimes p)_{jk}+C_{lijk}p_lV_{ijk}.
\end{equation}
\end{lemma}

\begin{proof}
Let $R=\Id-2p\otimes p$ be reflection in $p^\perp$.  Since $A$ is
symmetric and $Ap=\lambda p$, it preserves $p^\perp$, and hence
$RAR^T=A$.  Also $Rp=-p$.  Orthogonal equivariance and the even parity
in $p$ of $U,K$, and therefore of $G$, give
\[
 R^{\otimes4}G(A,p,0)=G(RAR^T,Rp,0)
                   =G(A,-p,0)=G(A,p,0).
\]
For the three-tensor $F_{ijk}=p_lG_{lijk}$, contraction commutes with
orthogonal transformations, so
\[
 R^{\otimes3}F=(Rp)_l(R^{\otimes4}G)_{lijk}=-F.
\]
Choose any orthonormal basis with its last vector $p$.  A component of a
three-tensor with $k$ normal slots is multiplied by $(-1)^k$ under
$R^{\otimes3}$.  It follows that $F$ has zero components whenever the
number of normal slots is even.  This assertion concerns the full
unsymmetrized tensor $F$; no permutation of its slots, and no extra
symmetry of $U$ or $K$, has been used.

For tangent indices $\alpha,\beta$, the boundary jet identities give
\[
 \frac{B_{N\alpha\beta}}{w^2}
 =\lambda A_{\alpha\beta}-(A^2)_{\alpha\beta}
 =(S_0)_{\alpha\beta},\qquad
 \frac{B_{NNN}}{w^2}=1.
\]
Because $S_0p=0$, these are precisely the components of $V$ with one or
three normal slots, including all their permutations.  Thus the part
of $B/w^2$ with an odd number of normal slots equals $V$; $V$ itself
has no even-normal components.  The even-normal components of $B$,
which are not determined by $A,p,g,w$, have
zero contraction with $F$.  This proves the first equality in
\eqref{flux:boundary-contraction}.

Finally, $S_0p=0$ and $|p|=1$ imply
\[
 p_iV_{ijk}=(S_0)_{jk}+p_jp_k.
\]
Splitting $G$ according to \eqref{flux:CG} proves the second equality.
The argument used only the orthogonal splitting
$\R^5=p^\perp\oplus\R p$.  It therefore applies without any assumption
on multiplicities, signs, or nonvanishing of the tangential eigenvalues
of $A$.
\end{proof}

\begin{proposition}\label{flux:boundary-comparison}
On the regular side boundary,
\begin{equation}\label{flux:signed-boundary-identity}
 N\cdot Z+Hh^2=w^{2a+1}\mathcal S
   \geq\delta_*w^{2a+1}>0,
 \qquad
 \delta_* =\frac{6984137359}{31250000000},
\end{equation}
where $\mathcal S$ is the boundary polynomial in
\eqref{cert:boundary}.  In particular $N\cdot Z+Hh^2\geq0$.
\end{proposition}

\begin{proof}
Proposition~\ref{cert:universal} gives $P=0$ at the boundary, where $g=0$.
Since $X=B/w$, Lemma~\ref{flux:boundary-parity} gives
\begin{align*}
 N\cdot Z
 &=w^{2a+1}\left(
       U_{jk}(S_0+p\otimes p)_{jk}
       +C_{lijk}p_lV_{ijk}+p_lM_l\right),\\
 Hh^2&=-\lambda w^{2a+1}f.
\end{align*}
Their sum is exactly $w^{2a+1}\mathcal S$.  The boundary data satisfy
all hypotheses of \eqref{cert:boundary}, proving the stated lower
bound.  The sign is expressed using the inward normal $N$; the outward
normal used by the divergence theorem is $-N$.

Only the boundary value of $P$ vanishes.  Its derivative term
$D_lP_l$ remains in \eqref{flux:divergence-identity}.  For example, a
monomial $P_l=gQ_l$ has $D_iP_l=p_iQ_l+gD_iQ_l$ at $z=1$, so its
derivative need not vanish when $g=0$.
\end{proof}

\subsection{Uniform estimates for the cutoff argument}

\begin{lemma}\label{flux:field-bounds}
There is a finite constant $C$, depending only on the coefficient
polynomials, the dimension, $a$, and $\epsilon_*$, such that on
$\Omega_+$
\begin{equation}\label{flux:uniform-field-bounds}
 |h|\leq Cw^a,\qquad
 |\nabla h|\leq Cw^a(1+w+|X|),\qquad
 |Z|\leq Cw^{2a}(1+w+|X|).
\end{equation}
The constant is independent of $w$, $|X|$, and all subsequent cutoff
parameters.
\end{lemma}

\begin{proof}
The coefficient domain
\[
 \{(A,p,g):A\in\STF^2(\R^5),\ |A|=1,\ |p|^2+g^2\leq1\}
\]
is compact.  The coefficient polynomials and their first derivatives
with respect to their coefficient variables are therefore uniformly
bounded there.  Also $|q|\leq|X|$ by Cauchy--Schwarz.  The rules
\eqref{flux:first-derivatives} consequently imply
\[
 |D f|+|q|\,|f|\leq C(1+w+|X|).
\]
Use \eqref{flux:f-positive} and
\eqref{flux:square-root-derivative} for the estimates on $h$ and
$\nabla h$, and use \eqref{flux:fields} and the boundedness of $G,M,P$
for the estimate on $Z$.
\end{proof}
